\Article{35}
Das folgende Problem verdient eine gesonderte Behandlung, sowohl wegen seines praktischen Nutzens als auch wegen der Eleganz seiner Lösung.

\textit{Die Änderungen der plausibelsten Werte der Unbekannten ermitteln, die durch das Hinzufügen einer neuen Gleichung entstehen, sowie die Gewichte der neuen Bestimmungen.}

Wir behalten die zuvor verwendeten Bezeichnungen bei, sodass die ursprünglichen, auf das Gewicht \(=1\) reduzierten Gleichungen diese sind: \(v=0,\allowbreak v'=0,\allowbreak v''=0\) etc.; für die unbestimmte Sum\SourceLeaf{048}%
me schreiben wir \(vv+v'v'+v''v''\quad\text{etc.}=\Omega\); ferner seien \(\xi,\eta,\zeta\) etc. die partiellen Differentialquotienten
\[
\frac{d\Omega}{2dx},\quad\frac{d\Omega}{2dy},\quad\frac{d\Omega}{2dz}\quad\text{etc.}
\]
und schließlich ergebe sich aus der unbestimmten Elimination
\[
\left.\begin{aligned}
x&=A+[\alpha\alpha]\xi+[\alpha\beta]\eta+[\alpha\gamma]\zeta+\text{etc.}\\
y&=B+[\alpha\beta]\xi+[\beta\beta]\eta+[\beta\gamma]\zeta+\text{etc.}\\
z&=C+[\alpha\gamma]\xi+[\beta\gamma]\eta+[\gamma\gamma]\zeta+\text{etc.}
\end{aligned}\right\}\ (\mathrm{I})
\]
Nehmen wir nun an, es komme eine neue Gleichung \(v^*=0\) hinzu (die näherungsweise richtig ist und das Gewicht \(=1\) hat), und untersuchen wir, wie stark sich dadurch sowohl die plausibelsten Werte der Unbekannten \(A,B,C\) etc. als auch die Koeffizienten \([\alpha\alpha],[\alpha\beta]\) etc. ändern.

Setzen wir \(\Omega+v^*v^*=\Omega^*,\allowbreak\frac{d\Omega^*}{2dx}=\xi^*,\allowbreak\frac{d\Omega^*}{2dy}=\eta^*,\allowbreak\frac{d\Omega^*}{2dz}=\zeta^*\) etc. und nehmen wir an, dass sich daraus durch Elimination ergibt
\[
x=A^*+[\alpha\alpha^*]\xi^*+[\alpha\beta^*]\eta^*+[\alpha\gamma^*]\zeta^*\quad\text{etc.}
\]
Schließlich sei
\[
v^*=fx+gy+hz+\text{etc.}+k
\]
sodass beim Einsetzen der Werte aus (I) für \(x,y,z\) etc. folgt
\[
v^*=F\xi+G\eta+H\zeta+\text{etc.}+K
\]
und setzen wir
\[
Ff+Gg+Hh+\text{etc.}=\omega
\]
Offenbar ist \(K\) der plausibelste Wert der Funktion \(v^*\), soweit er sich aus den ursprünglichen Gleichungen ergibt, ohne Rücksicht auf den Wert \(0\), den die zusätzliche Beobachtung geliefert hat, und \(\frac1\omega\) ist das Gewicht dieser Bestimmung.

Nun haben wir
\[
\xi^*=\xi+fv^*,\quad\eta^*=\eta+gv^*,\quad\zeta^*=\zeta+hv^*\quad\text{etc.}
\]
und daher
\[
\begin{aligned}
&F\xi^*+G\eta^*+H\zeta^*+\text{etc.}+K\\
&\qquad=v^*(1+Ff+Gg+Hh+\text{etc.})
\end{aligned}
\]
oder
\[
v^*=\frac{F\xi^*+G\eta^*+H\zeta^*+\text{etc.}+K}{1+\omega}
\]
\SourceLeaf{049}
\noindent Ebenso ergibt sich
\[
\begin{aligned}
x&=A+[\alpha\alpha]\xi^*+[\alpha\beta]\eta^*+[\alpha\gamma]\zeta^*+\text{etc.}\\
&\qquad-v^*(f[\alpha\alpha]+g[\alpha\beta]+h[\alpha\gamma]+\text{etc.})\\
&=A+[\alpha\alpha]\xi^*+[\alpha\beta]\eta^*+[\alpha\gamma]\zeta^*+\text{etc.}-Fv^*\\
&=A+[\alpha\alpha]\xi^*+[\alpha\beta]\eta^*+[\alpha\gamma]\zeta^*+\text{etc.}\\
&\qquad-\frac{F}{1+\omega}(F\xi^*+G\eta^*+H\zeta^*+\text{etc.}+K)
\end{aligned}
\]
Daraus schließen wir also
\[
A^*=A-\frac{FK}{1+\omega}
\]
was der plausibelste Wert von \(x\) aus \textit{allen} Beobachtungen ist;
\[
[\alpha\alpha^*]=[\alpha\alpha]-\frac{FF}{1+\omega}
\]
und das Gewicht dieser Bestimmung ist daher
\[
=\frac1{[\alpha\alpha]-\dfrac{FF}{1+\omega}}
\]
Auf ganz dieselbe Weise findet man als plausibelsten Wert von \(y\), der auf \textit{allen} Beobachtungen beruht,
\[
B^*=B-\frac{GK}{1+\omega}
\]
und als Gewicht dieser Bestimmung
\[
=\frac1{[\beta\beta]-\dfrac{GG}{1+\omega}}
\]
und so fort. Q.\ E.\ I.

Dieser Lösung seien einige Bemerkungen hinzugefügt.

I. Setzt man diese neuen Werte \(A^*,B^*,C^*\) etc. ein, so erhält die Funktion \(v^*\) den plausibelsten Wert
\[
K-\frac{K}{1+\omega}(Ff+Gg+Hh+\text{etc.})=\frac{K}{1+\omega}
\]
Da ferner unbestimmt gilt
\[
v^*=\frac{F}{1+\omega}.\xi^*+\frac{G}{1+\omega}.\eta^*+\frac{H}{1+\omega}.\zeta^*+\text{etc.}+\frac{K}{1+\omega}
\]
\SourceLeaf{050}%
ergibt sich das Gewicht dieser Bestimmung nach den Grundsätzen von Art.\ 29 als
\[
=\frac{1+\omega}{Ff+Gg+Hh+\text{etc.}}=\frac1\omega+1
\]
Dasselbe folgt unmittelbar aus der Anwendung der am Ende von Art.\ 21 angegebenen Regel: Das System der ursprünglichen Gleichungen hatte die Bestimmung \(v^*=K\) mit dem Gewicht \(=\frac1\omega\) geliefert; danach gab die neue Beobachtung eine andere, von jener unabhängige Bestimmung \(v^*=0\) mit dem Gewicht \(=1\); ihre Verbindung ergibt die Bestimmung \(v^*=\frac{K}{1+\omega}\) mit dem Gewicht \(=\frac1\omega+1\).

II. Da für \(x=A^*,\allowbreak y=B^*,\allowbreak z=C^*\) etc. auch \(\xi^*=0,\allowbreak\eta^*=0,\allowbreak\zeta^*=0\) etc. sein muss, folgt ferner für dieselben Werte
\[
\xi=-\frac{fK}{1+\omega},\quad\eta=-\frac{gK}{1+\omega},\quad\zeta=-\frac{hK}{1+\omega}\quad\text{etc.}
\]
und außerdem wegen der unbestimmt geltenden Gleichung
\(\Omega=\xi(x-A)+\eta(y-B)+\zeta(z-C)+\text{etc.}+M\),
\[
\begin{aligned}
\Omega&=\frac{KK}{(1+\omega)^2}(Ff+Gg+Hh+\text{etc.})+M\\
&=M+\frac{\omega KK}{(1+\omega)^2}
\end{aligned}
\]
schließlich wegen der unbestimmt geltenden Gleichung
\(\Omega^*=\Omega+v^*v^*\),
\[
\Omega^*=M+\frac{\omega KK}{(1+\omega)^2}+\frac{KK}{(1+\omega)^2}=M+\frac{KK}{1+\omega}
\]
III. Vergleichen wir dies mit dem in Art.\ 30 Gelehrten, so bemerken wir, dass die Funktion \(\Omega\) hier den kleinsten Wert annimmt, den sie für den bestimmten Wert der Funktion \(v^*=\frac{K}{1+\omega}\) annehmen kann.

\Article{36}
Für ein anderes, mit dem vorhergehenden verwandtes Problem, nämlich

\textit{Die Änderungen der plausibelsten Werte der Unbekannten untersuchen, die aus einer Änderung des Gewichts einer der ursprünglichen Beobachtungen entstehen, sowie die Gewichte der neuen Bestimmungen}

\noindent geben wir hier nur die Lösung an und lassen der Kürze wegen den Beweis weg, der sich nach dem Muster des vorhergehenden Artikels leicht durchführen lässt.

Nehmen wir an, erst nach Abschluss der Rechnung bemerke man, dass einer der Beobachtungen ein zu kleines oder zu großes Gewicht zugewiesen wurde; dass z.\ B. der ersten Beobachtung, die \(V=L\) ergab, anstelle des in der Rechnung verwendeten Gewichts \(p\) richtiger das Gewicht \(=p^*\) zuzuweisen ist. Dann braucht man nicht die ganze Rechnung zu wiederholen, sondern kann die Korrekturen bequemer mit den folgenden Formeln berechnen.

\SourceLeaf{051}%
Die korrigierten plausibelsten Werte der Unbekannten sind diese:
\[
\begin{aligned}
x&=A-\frac{(p^*-p)\alpha\lambda}{p+(p^*-p)(a\alpha+b\beta+c\gamma+\text{etc.})}\\[7pt]
y&=B-\frac{(p^*-p)\beta\lambda}{p+(p^*-p)(a\alpha+b\beta+c\gamma+\text{etc.})}\\[7pt]
z&=C-\frac{(p^*-p)\gamma\lambda}{p+(p^*-p)(a\alpha+b\beta+c\gamma+\text{etc.})}
\end{aligned}
\]
etc.; die Gewichte dieser Bestimmungen findet man, indem man die Einheit jeweils durch die folgenden Ausdrücke dividiert:
\[
\begin{aligned}
&[\alpha\alpha]-\frac{(p^*-p)\alpha\alpha}{p+(p^*-p)(a\alpha+b\beta+c\gamma+\text{etc.})}\\[7pt]
&[\beta\beta]-\frac{(p^*-p)\beta\beta}{p+(p^*-p)(a\alpha+b\beta+c\gamma+\text{etc.})}\\[7pt]
&[\gamma\gamma]-\frac{(p^*-p)\gamma\gamma}{p+(p^*-p)(a\alpha+b\beta+c\gamma+\text{etc.})}\quad\text{etc.}
\end{aligned}
\]
Diese Lösung umfasst zugleich den Fall, dass man nach Abschluss der Rechnung erkennt, eine der Beobachtungen hätte gänzlich verworfen werden müssen; denn dies ist dasselbe, wie \(p^*=0\) zu setzen. Ebenso bezieht sich der Wert \(p^*=\infty\) auf den Fall, dass die Gleichung \(V=L\), die in der Rechnung als näherungsweise richtig behandelt wurde, tatsächlich absolute Genauigkeit besitzt.

Kommen im Übrigen zu den Gleichungen, auf denen die Rechnung beruhte, \textit{mehrere} neue hinzu, oder erkennt man, dass \textit{mehreren} von ihnen falsche Gewichte zugewiesen wurden, so würde die Berechnung der Korrekturen zu kompliziert; in einem solchen Fall wird es daher besser sein, die Rechnung von Grund auf neu auszuführen.

\clearpage
